\section{审判}
\begin{enumerate}
  \item 审判 = 审理 + 裁判。
  \item 特征。
        \begin{enumerate}
          \item 被动性：\zd{不告不理}；独立性；中立性；职权性；程序性；
          \item \zd{亲历性}：自始自终参与，充分听取控辩双方意见；公开性；终局性。
        \end{enumerate}
  \item 基本的审判程序：一审、二审、复核（死刑、法定刑以下）、审判监督程序。有先后，可跳跃，可回转，可夭折。
  \item 模式。控辩审三方地位、关系。同基本范畴里的刑事构造。
  \item 原则。
        \begin{enumerate}
          \item 审判公开。国家秘密、个人隐私、开庭时未满 18 \yjt 绝对；商业秘密 \yjt 相对不公开。
          \item 直接言词。亲自面对，口头陈述。
          \item 集中审理。人不换事不断。
          \item 辩论。控辩双方及当事人就\zd{证据、事实、程序、法律适用}吵架，贯穿始终。
        \end{enumerate}
  \item 审级制度。两审终审。
  \item 审判组织。
        \begin{enumerate}
          \item 独任制。基层简速 3 下独。简易 3 年以下。
          \item 合议制。
                \begin{enumerate}
                  \item 组成方式。
                        \begin{enumerate}
                          \item \zd{地方一审可以陪}；基中 3 高到 57；法陪统统 3 或 7；
                          \item 二审法官 3 或 5；复核法官仅 3 人；
                          \item 再审具体看几审；重审另组合议庭；
                        \end{enumerate}
                  \item 组成原则。随机组成。
                  \item 活动原则。
                  \item 免责事由。\zd{无过才无责}。法律、事实、证据认识偏差。
                \end{enumerate}
          \item 审判委员会。\zd{抗}诉、\zd{死}刑、\zd{要再审}必上审委会；评议后、不留名、非必经、需执行。
        \end{enumerate}
  \item 人民陪审员。
        \begin{enumerate}
          \item 任职条件。28 岁；高中以上文化程度。公道正派。
          \item 适用范围。地方一审可以陪。
                \begin{enumerate}
                  \item 依职权：
                        \begin{enumerate}
                          \item \zd{群体、公共}利益；\zd{影响大；案情复杂或其他}。
                          \item \zd{征}地拆迁、生态\zd{环}境、\zd{药}品\zd{食}品、10 年以上/死刑、公\zd{益}诉讼。\zd{甄环要 10 亿，法陪 7 人合议庭}。
                        \end{enumerate}
                  \item 依申请。
                \end{enumerate}
          \item 选任程序。略。
          \item 权利义务。3 人事实 + 法律；7 人仅事实。
        \end{enumerate}
\end{enumerate}
